\documentclass[11pt,a4paper]{article}
% Preambulo comun de los reportes de Tecnologias Celulares.
% Cada main.tex hace: \documentclass ... % Preambulo comun de los reportes de Tecnologias Celulares.
% Cada main.tex hace: \documentclass ... % Preambulo comun de los reportes de Tecnologias Celulares.
% Cada main.tex hace: \documentclass ... \input{../comun/preambulo.tex}
\usepackage[utf8]{inputenc}
\usepackage[T1]{fontenc}
\usepackage[spanish,mexico]{babel}
\usepackage{graphicx}
\usepackage[margin=2.5cm]{geometry}
\usepackage{amsmath}
\usepackage{amssymb}
\usepackage{booktabs}
\usepackage{tabularx}
\usepackage{array}
\usepackage{float}
\usepackage{setspace}
\usepackage{parskip}
\usepackage{fancyhdr}
\usepackage{textcomp}
\usepackage{xcolor}
\usepackage{caption}
\usepackage{fvextra}
\usepackage[hidelinks]{hyperref}

\onehalfspacing
\setlength{\parskip}{6pt}
\sloppy
\setlength{\emergencystretch}{3em}
\setlength{\headheight}{40pt}
\captionsetup{font=small,labelfont=bf}
\setlength{\tabcolsep}{6pt}

\pagestyle{fancy}
\fancyhf{}
\fancyhead[L]{\includegraphics[width=1cm]{/home/yovick/repos/UAZ-IEI/Assets/UIE.png}}
\fancyhead[R]{\small Tecnologías Celulares}
\fancyfoot[C]{\thepage}
\renewcommand{\headrulewidth}{0.4pt}

\DefineVerbatimEnvironment{codigo}{Verbatim}{%
  breaklines=true,breakanywhere=true,fontsize=\small,frame=single,
  framesep=3pt,rulecolor=\color{gray!50}}

% Figura centrada con pie: \figura[ancho]{archivo}{leyenda}
\newcommand{\figura}[3][0.85]{%
  \begin{center}
    \includegraphics[width=#1\linewidth]{#2}\\[2pt]
    {\small #3}
  \end{center}}

% Portada: \portada{numero}{titulo de la practica}
\newcommand{\portada}[2]{%
\begin{titlepage}
  \thispagestyle{empty}
  \begin{center}
    \vspace*{1cm}
    \includegraphics[width=4.5cm]{/home/yovick/repos/UAZ-IEI/Assets/uaz.png}
    \vspace{0.8cm}

    {\large Universidad Autónoma de Zacatecas}\\[4pt]
    {\large Unidad Académica de Ingeniería Eléctrica}\\[4pt]
    {\normalsize Ingeniería Electrónica Industrial con especialidad en Telecomunicaciones}

    \vspace{1.8cm}
    {\LARGE\bfseries Práctica #1}\\[10pt]
    {\Large\bfseries #2}
  \end{center}

  \vspace{2.2cm}
  \begin{center}
    \begin{tabular}{rl}
      \textbf{Materia:} & Tecnologías Celulares\\
      \textbf{Docente:} & Dr. Víktor Iván Rodríguez Abdalá\\
      \textbf{Alumno:}  & Arnold Yovick Ríos Zaldívar\\
      \textbf{Fecha:}   & Noviembre de 2026\\
    \end{tabular}
  \end{center}
  \vfill
\end{titlepage}}

\usepackage[utf8]{inputenc}
\usepackage[T1]{fontenc}
\usepackage[spanish,mexico]{babel}
\usepackage{graphicx}
\usepackage[margin=2.5cm]{geometry}
\usepackage{amsmath}
\usepackage{amssymb}
\usepackage{booktabs}
\usepackage{tabularx}
\usepackage{array}
\usepackage{float}
\usepackage{setspace}
\usepackage{parskip}
\usepackage{fancyhdr}
\usepackage{textcomp}
\usepackage{xcolor}
\usepackage{caption}
\usepackage{fvextra}
\usepackage[hidelinks]{hyperref}

\onehalfspacing
\setlength{\parskip}{6pt}
\sloppy
\setlength{\emergencystretch}{3em}
\setlength{\headheight}{40pt}
\captionsetup{font=small,labelfont=bf}
\setlength{\tabcolsep}{6pt}

\pagestyle{fancy}
\fancyhf{}
\fancyhead[L]{\includegraphics[width=1cm]{/home/yovick/repos/UAZ-IEI/Assets/UIE.png}}
\fancyhead[R]{\small Tecnologías Celulares}
\fancyfoot[C]{\thepage}
\renewcommand{\headrulewidth}{0.4pt}

\DefineVerbatimEnvironment{codigo}{Verbatim}{%
  breaklines=true,breakanywhere=true,fontsize=\small,frame=single,
  framesep=3pt,rulecolor=\color{gray!50}}

% Figura centrada con pie: \figura[ancho]{archivo}{leyenda}
\newcommand{\figura}[3][0.85]{%
  \begin{center}
    \includegraphics[width=#1\linewidth]{#2}\\[2pt]
    {\small #3}
  \end{center}}

% Portada: \portada{numero}{titulo de la practica}
\newcommand{\portada}[2]{%
\begin{titlepage}
  \thispagestyle{empty}
  \begin{center}
    \vspace*{1cm}
    \includegraphics[width=4.5cm]{/home/yovick/repos/UAZ-IEI/Assets/uaz.png}
    \vspace{0.8cm}

    {\large Universidad Autónoma de Zacatecas}\\[4pt]
    {\large Unidad Académica de Ingeniería Eléctrica}\\[4pt]
    {\normalsize Ingeniería Electrónica Industrial con especialidad en Telecomunicaciones}

    \vspace{1.8cm}
    {\LARGE\bfseries Práctica #1}\\[10pt]
    {\Large\bfseries #2}
  \end{center}

  \vspace{2.2cm}
  \begin{center}
    \begin{tabular}{rl}
      \textbf{Materia:} & Tecnologías Celulares\\
      \textbf{Docente:} & Dr. Víktor Iván Rodríguez Abdalá\\
      \textbf{Alumno:}  & Arnold Yovick Ríos Zaldívar\\
      \textbf{Fecha:}   & Noviembre de 2026\\
    \end{tabular}
  \end{center}
  \vfill
\end{titlepage}}

\usepackage[utf8]{inputenc}
\usepackage[T1]{fontenc}
\usepackage[spanish,mexico]{babel}
\usepackage{graphicx}
\usepackage[margin=2.5cm]{geometry}
\usepackage{amsmath}
\usepackage{amssymb}
\usepackage{booktabs}
\usepackage{tabularx}
\usepackage{array}
\usepackage{float}
\usepackage{setspace}
\usepackage{parskip}
\usepackage{fancyhdr}
\usepackage{textcomp}
\usepackage{xcolor}
\usepackage{caption}
\usepackage{fvextra}
\usepackage[hidelinks]{hyperref}

\onehalfspacing
\setlength{\parskip}{6pt}
\sloppy
\setlength{\emergencystretch}{3em}
\setlength{\headheight}{40pt}
\captionsetup{font=small,labelfont=bf}
\setlength{\tabcolsep}{6pt}

\pagestyle{fancy}
\fancyhf{}
\fancyhead[L]{\includegraphics[width=1cm]{/home/yovick/repos/UAZ-IEI/Assets/UIE.png}}
\fancyhead[R]{\small Tecnologías Celulares}
\fancyfoot[C]{\thepage}
\renewcommand{\headrulewidth}{0.4pt}

\DefineVerbatimEnvironment{codigo}{Verbatim}{%
  breaklines=true,breakanywhere=true,fontsize=\small,frame=single,
  framesep=3pt,rulecolor=\color{gray!50}}

% Figura centrada con pie: \figura[ancho]{archivo}{leyenda}
\newcommand{\figura}[3][0.85]{%
  \begin{center}
    \includegraphics[width=#1\linewidth]{#2}\\[2pt]
    {\small #3}
  \end{center}}

% Portada: \portada{numero}{titulo de la practica}
\newcommand{\portada}[2]{%
\begin{titlepage}
  \thispagestyle{empty}
  \begin{center}
    \vspace*{1cm}
    \includegraphics[width=4.5cm]{/home/yovick/repos/UAZ-IEI/Assets/uaz.png}
    \vspace{0.8cm}

    {\large Universidad Autónoma de Zacatecas}\\[4pt]
    {\large Unidad Académica de Ingeniería Eléctrica}\\[4pt]
    {\normalsize Ingeniería Electrónica Industrial con especialidad en Telecomunicaciones}

    \vspace{1.8cm}
    {\LARGE\bfseries Práctica #1}\\[10pt]
    {\Large\bfseries #2}
  \end{center}

  \vspace{2.2cm}
  \begin{center}
    \begin{tabular}{rl}
      \textbf{Materia:} & Tecnologías Celulares\\
      \textbf{Docente:} & Dr. Víktor Iván Rodríguez Abdalá\\
      \textbf{Alumno:}  & Arnold Yovick Ríos Zaldívar\\
      \textbf{Fecha:}   & Noviembre de 2026\\
    \end{tabular}
  \end{center}
  \vfill
\end{titlepage}}


\begin{document}

\portada{5}{Generación de una señal OFDM}

\textbf{Script:} \texttt{practica5.py} (Python con NumPy/SciPy)

\section{Objetivo}

Implementar un transmisor OFDM basado en la IFFT, generar símbolos QPSK sobre 64
subportadoras y verificar la ortogonalidad entre subportadoras y las características
temporales y espectrales de la señal resultante.

\section{Fundamento teórico}

Una señal OFDM es la suma de $N$ subportadoras ortogonales moduladas por símbolos
complejos:
\begin{equation}
  s(t) = \sum_{k=0}^{N-1} X_k\,e^{j2\pi k t/T},
  \qquad 0 \le t < T.
\end{equation}
Muestreando en $t_n=nT/N$ queda
\begin{equation}
  x[n] = \frac{1}{\sqrt{N}}\sum_{k=0}^{N-1} X_k\,e^{j2\pi kn/N},
\end{equation}
que es exactamente la transformada discreta inversa de Fourier multiplicada por
$\sqrt{N}$. La IFFT permite generar en $O(N\log N)$ lo que en un transmisor analógico
requeriría $N$ osciladores.

La ortogonalidad se comprueba con el producto interno
\begin{equation}
  \langle s_k, s_m\rangle = \sum_{n=0}^{N-1} e^{j2\pi kn/N}\,e^{-j2\pi mn/N}
  = N\,\delta_{km},
\end{equation}
de modo que cada subportadora se puede separar en el receptor sin interferencia de las
demás, aunque sus espectros se superpongan.

\section{Desarrollo experimental}

\subsection{Parámetros}

\begin{tabularx}{\linewidth}{@{}*{2}{>{\raggedright\arraybackslash}X}@{}}
  \toprule
  Parámetro & Valor \\
  \midrule
  Bits generados & 1024 \\
  Modulación & QPSK con mapeo Gray \\
  Símbolos QPSK & 512 \\
  Subportadoras ($N_{\mathrm{FFT}}$) & 64 \\
  Símbolos OFDM & 8 \\
  Muestras transmitidas & 512 \\
  \bottomrule
\end{tabularx}

\subsection{Procedimiento}

\begin{enumerate}
  \item Los 1024 bits se agrupan en pares y se convierten en 512 símbolos QPSK
        $(\pm1\pm j)/\sqrt{2}$.
  \item Los símbolos se ordenan en una matriz de $64\times8$; cada columna es un símbolo
        OFDM.
  \item Se aplica la IFFT por columnas y se serializa la señal en orden de columnas.
  \item Se grafican la señal temporal, la constelación, la PSD, las magnitudes de las
        subportadoras y los productos internos.
\end{enumerate}

\section{Código}

\begin{codigo}
def modular_qpsk(bits):
    b = np.asarray(bits).reshape(-1, 2)
    return ((2 * b[:, 0] - 1) + 1j * (2 * b[:, 1] - 1)) / np.sqrt(2)

def generar_ofdm(simbolos, nfft=64):
    x = np.fft.ifft(simbolos.reshape(-1, nfft).T, axis=0) * np.sqrt(nfft)
    return x.flatten(order="F"), x
\end{codigo}

\section{Resultados}

\subsection{Señal OFDM en el tiempo}

\figura{figs/senal_ofdm.png}{Parte real de los primeros cuatro símbolos OFDM. Las líneas punteadas marcan los límites entre símbolos.}

La señal tiene una apariencia ruidosa, con picos que superan varias veces el valor eficaz.
Esa suma constructiva de 64 subportadoras es la causa de la relación de potencia pico a
promedio (PAPR) que se estudia en la práctica 8; en esta corrida resultó de 8.21 dB.

\subsection{Constelación QPSK}

\figura[0.55]{figs/constelacion.png}{Constelación de los 512 símbolos QPSK transmitidos.}

Los símbolos se ubican exactamente en las cuatro esquinas $(\pm1,\pm1)/\sqrt{2}$, como
corresponde a una modulación QPSK sin ruido.

\subsection{Espectro}

\figura{figs/espectro.png}{Densidad espectral de potencia de la señal OFDM.}

El espectro es aproximadamente plano dentro de la banda ocupada, con las ondulaciones
propias de la suma finita de subportadoras y los lóbulos de las transiciones entre
símbolos. Ocupa todo el ancho de banda de muestreo, que es la firma de una señal
multiportadora.

\subsection{Subportadoras}

\figura{figs/subportadoras.png}{Magnitud de las 64 subportadoras del primer símbolo OFDM.}

\subsection{Ortogonalidad}

\figura{figs/ortogonalidad.png}{Dos subportadoras consecutivas (arriba) y producto interno de cada subportadora con la subportadora 5 (abajo).}

El producto interno $\langle s_5,s_6\rangle$ resultó de \textbf{7.41$\times10^{-15}$} y el
máximo producto fuera de la diagonal, \textbf{1.06$\times10^{-13}$}. Ambos son del orden
del error de redondeo de punto flotante, es decir, cero en la práctica: las subportadoras
son ortogonales.

\section{Análisis de resultados}

\begin{itemize}
  \item \textbf{¿Por qué las subportadoras no interfieren entre sí?} Porque sus funciones
        base son ortogonales en el intervalo del símbolo: el producto interno entre dos
        subportadoras distintas es cero. Cada una tiene un número entero de ciclos en $T$,
        de modo que los lóbulos de una caen en los nulos de las demás.
  \item \textbf{¿Cuál es la función de la IFFT?} Convertir los símbolos en el dominio de la
        frecuencia (una muestra por subportadora) en la forma de onda temporal que se
        transmite. Es la implementación eficiente de la suma de subportadoras.
  \item \textbf{¿Qué diferencias existen entre OFDM y una transmisión monocanal?} En
        monocanal un solo flujo ocupa todo el ancho de banda con un símbolo largo; en OFDM
        muchos flujos de baja tasa se transmiten en paralelo, cada uno en una subportadora
        angosta. Eso alarga la duración del símbolo y reduce la sensibilidad a la
        dispersión del canal.
  \item \textbf{¿Cómo cambia el espectro al aumentar el número de subportadoras?} La banda
        total depende de la separación entre subportadoras ($\Delta f = F_s/N$), no del
        número: al aumentar $N$ con la misma $F_s$, cada subportadora se hace más angosta
        y hay más, pero la ocupación total es la misma. Lo que cambia es la densidad de
        subportadoras y la duración del símbolo.
\end{itemize}

\section{Cuestionario}

\begin{enumerate}
  \item \textbf{¿Qué es OFDM?} Multiplexación por división de frecuencia ortogonal: una
        técnica multiportadora que divide el flujo de datos en muchos flujos paralelos
        transmitidos simultáneamente en subportadoras ortogonales.
  \item \textbf{¿Por qué se utilizan subportadoras ortogonales?} Para que los espectros se
        puedan superponer sin interferencia mutua y el receptor pueda separarlas con una
        FFT, aprovechando el espectro de forma más eficiente que con canales separados por
        bandas de guarda.
  \item \textbf{¿Cuál es la función de la IFFT?} Generar la señal temporal a partir de los
        símbolos en frecuencia. Es la operación inversa que el receptor deshace con la FFT.
  \item \textbf{¿Qué ventajas ofrece OFDM en canales selectivos en frecuencia?} Convierte
        un canal selectivo en muchos subcanales aproximadamente planos, de modo que la
        ecualización se reduce a un factor complejo por subportadora. Además, el símbolo
        largo tolera mejor la dispersión temporal.
  \item \textbf{¿Qué tecnologías utilizan OFDM actualmente?} LTE y 5G NR (OFDMA/SC-FDMA en
        el enlace descendente/ascendente), Wi-Fi (802.11a/g/n/ac/ax), DVB-T, WiMAX, ADSL y
        PLC.
\end{enumerate}

\section{Conclusiones}

Se generó una señal OFDM de 8 símbolos con 64 subportadoras QPSK a partir de 1024 bits.
La IFFT produjo una señal temporal con PAPR de 8.21 dB y un espectro plano dentro de la
banda. La verificación numérica de ortogonalidad dio $7.41\times10^{-15}$ para el producto
entre las subportadoras 5 y 6 y $1.06\times10^{-13}$ como máximo fuera de la diagonal, lo
que confirma que las subportadoras son ortogonales y que la IFFT es equivalente a la suma
de moduladores independientes.

\end{document}
